% PAPER_A_METHOD_CONTRACT_CURRENT.json paper_a_method_20260730_v5 SHA-256 bde2e7e20cb00fa4f67b377112fe6534e27e7938c34fb4f63b7987fd3c142e2b
% Paper A: Strong Global is the EB external-quality baseline.
\section{Strong Global 与 Support-aware Full}

\subsection{Global 候选审计}

Global-LOO、Global-GT 和 Global-GT-gated-LOO 仍可在 cross-fitted replay 中比较，
报告排名重叠、discrepancy worker 和稳定性；这些只是审计和敏感性分析，不能替代正式
Strong Global，也不能由 V1 outcome 选择。

\subsection{Strong Global}

\subsubsection{Eligibility}

\begin{verbatim}
process valid
independence valid
GT support sufficient
building/task coverage sufficient
no recurrent serious structural failure
Q_GT_EB meets the production quality floor
\end{verbatim}

\subsubsection{正式分数与安全门}

\[
S^G_u=z\!\left(\widehat Q^{GT,EB}_u\right).
\]

其中 \(z(\cdot)\) 的均值和标准差只在冻结的、行政上具备资格且可估计
\(Q^{GT}\) 的 capability cohort 上计算，并在质量 floor 与 structural gate 之前完成。
通过质量 floor 或 structural gate 后不重新标准化；Strong Global 的正式排序不能借由
筛选后的资格变化改写尺度。

正式排序使用收缩后的外部正确性；LCB 只作为最低质量安全门、不确定性标记和保守敏感性
排序。同行指标不进入正式分数，仅用于 tie-break、discrepancy worker、聚合稳定性和敏感性。

\subsection{Support-aware Full}

\[
S^F_{u,t}=S^G_u
+I_{\mathrm{risk},t}B_u
+I_{f,t}H_{u,f}.
\]

其中 \(B_u\) 表示 ordinary/stress 风险韧性，\(H_{u,f}\) 表示一个已激活 failure-family
的条件能力。每任务最多激活一个 family，总调整有上限；证据有限时向群体均值收缩，
component 无支持时该项为零，任务超出 Calibration 支持范围时整体 fallback Strong Global。

\subsection{可路由 family 与激活}

正式 Full 只允许：

\begin{enumerate}
  \item undercoverage / underextension；
  \item adjacent-space overextension；
  \item corner / topology instability。
\end{enumerate}

\texttt{blind trust} 与 \texttt{correction failure} 仅作为 T1 机制结果；Scope 和
process-integrity 不进入 V1 条件加权。任务侧 family 必须唯一激活；激活含糊则关闭该项。

\subsection{支持、权重与 fallback}

跨阶段 component support 是预声明的方法合同。对每一个预定义 component，使用任务级
Calibration outcome 的 stage-stacked 回归或层级模型，包含 P1 worker component、stage、
P1\(\times\)stage interaction，并按 task 与 worker 聚类或设置相应随机效应。模型输出
方向、stage interaction、worker/task support 与收缩估计，并在 Main outcome 可见前冻结。
若 C2 出现清晰的方向反转、支持不足或不稳定收缩，该 component 不进入 Full。

只有当上述模型工件和独立 policy manifest 均有效时，component 才能通过 P1 integrity、
C1 predictive validation 和 C2-B confirmation，并具备足够 worker/task support、方向稳定、
收缩估计和标注前可激活性。权重与总调整上限不能为了增加政策差异而改变。

\begin{verbatim}
component unsupported        -> component adjustment = 0
family activation ambiguous  -> component adjustment = 0
d_cal_F out of support       -> overall fallback Strong Global
profile version incompatible -> overall fallback Strong Global
ranking unstable             -> overall fallback Strong Global
\end{verbatim}

整体 fallback 时，Full 的排序逐项等于 Strong Global。C2 closeout 后仅报告 activation、
fallback、推荐不同率、supported candidate count 和 capacity 后差异，不能用 V1 结果放宽门槛。

\subsection{Stage 3 formal manifest contract}

Strong Global 的 \(z\)-score cohort 在 quality/structural gate 前冻结为 administratively eligible 且
\(Q^{GT,EB}\) estimable 的 worker；cohort 少于两人、标准差为零、字段缺失或非有限值时 fail closed。
LCB 只作 quality floor、区间和 sensitivity，不参与正式主排序。Full manifest 必须显式绑定
component whitelist、worker/task support、weight、activation threshold/margin、symmetric cap、
component interval、profile version、fallback、formal minimum worker/cluster rule 与全部输入 SHA。
任何字段缺失、版本冲突、支持不足或 ranking endpoint 不稳定时，Full 整体排序精确回到 Strong Global。
V1 freeze manifest 还必须消费开放且 SHA 匹配的 Stage 3 gate；GT-blind aggregation 不读取
\(Q^{GT}\)、\(S^G\)、\(S^F\) 或 worker quality 作 medoid tie-break。

\subsection{静态画像与运行时调度分离}

Strong Global 的唯一静态顺序为
\(S_G\rightarrow R_{\mathrm{peer,stable}}\rightarrow R_{\mathrm{LOO,medoid}}\rightarrow\)
frozen random。availability 与 capacity 只在运行时 scheduler 过滤候选，不进入冻结画像或
质量排序。V1 formal candidate 只消费 \(S_G\)，旧 \texttt{global\_lcb} 字段必须拒绝。

\paragraph{Contract v5 clarification.}
Strong Global is frozen once, with order \(S_G\rightarrow R_{peer,stable}\rightarrow R_{LOO,medoid}\rightarrow\) frozen random; a tie-break layer is used only when every candidate in the current tie group is evaluable. An unsupported risk/family component, ambiguous family activation, or insufficient conditional support sets only that component to zero. Overall Full fallback is reserved for out-of-support tasks, profile-version conflict, or endpoint ranking instability.
